\section{Code ensembles and distance events}
\label{sec:constructions}

This section defines the objects sampled in the paper.  We use row vectors,
so a message $\x\in\F_p^k$ is encoded as $\x\G\in\F_p^n$ by a matrix
$\G\in\F_p^{k\times n}$.
Throughout the paper, $p$ denotes the order of the field.  It may be a prime
power; it does not necessarily denote the characteristic.

\subsection{Distance and the GV benchmark}

For $\y\in\F_p^n$, its support is
$\supp(\y):=\{i:y_i\ne0\}$, and its Hamming weight is
$\wt(\y):=|\supp(\y)|$.  Vectors of the same weight form a \emph{Hamming
shell}.  Define
\begin{equation}
 d_{\min}(\G):=
 \min_{\x\in\F_p^k\setminus\{0\}}\wt(\x\G).
 \label{eq:minimum-distance}
\end{equation}
This convention assigns distance zero to a noninjective generator.  If
$d_{\min}(\G)>0$, the matrix is injective and
\eqref{eq:minimum-distance} equals the usual minimum distance of its image.
The relative minimum distance is $d_{\min}(\G)/n$.

For an integer cutoff $L$, define the bad-sampling event
\begin{equation}
 \Bad_L(\G):=\{d_{\min}(\G)\leq L\}.
 \label{eq:bad-event}
\end{equation}
Every probability in this paper is over the stated random choices used to
sample $\G$.  A bound $\Pr[\Bad_L(\G)]\leq U$ shows that a sampled generator
has distance above $L$ except with probability at most $U$.

The code rate is $R:=k/n$.  We use natural logarithms and define
\begin{align}
 h(x)&:=-x\ln x-(1-x)\ln(1-x),\label{eq:binary-entropy}\\
 h_p(x)&:=h(x)+x\ln(p-1).\label{eq:p-ary-entropy}
\end{align}
The $p$-ary GV relative distance $\delta_{p,\mathrm{GV}}(R)$ is the solution
in $(0,(p-1)/p)$ to
\begin{equation}
 h_p(\delta_{p,\mathrm{GV}}(R))=(1-R)\ln p.
 \label{eq:gv-definition}
\end{equation}
This equation balances the number $p^k$ of messages against the exponential
volume of a radius-$\delta n$ Hamming ball.  We use it only as a distance
benchmark.
Here a radius-$L$ Hamming ball centered at zero is the set of vectors of
weight at most $L$; it contains
$\sum_{h=0}^{L}\binom nh(p-1)^h$ vectors.

\subsection{The common architecture}

Every construction has the form
\begin{equation}
 \G=\B\T,
 \label{eq:common-architecture}
\end{equation}
where $\B\in\F_p^{k\times n}$ is sparse and
$\T\in\F_p^{n\times n}$ is an invertible recursive map.  We call $\B$ the
expander matrix, following the EA and EC literature.  No spectral or
combinatorial expansion property is implicit in that name.

The rows of $\B$ correspond to message coordinates.  Its columns correspond
to the intermediate coordinates supplied to $\T$.  Thus the bipartite graph
of $\B$ has $k$ left vertices and $n$ right vertices.  A matrix entry is an
edge when it is nonzero.

\paragraph{Accumulator.}
Given $\mathbf u\in\F_p^n$, initialize $y_0:=0$ and compute
\begin{equation}
 y_t:=u_t+y_{t-1},\qquad t\in\{1,\ldots,n\}.
 \label{eq:accumulator-definition}
\end{equation}
The resulting linear map is the accumulator $\A$.  It is invertible because
$u_t=y_t-y_{t-1}$.  The EA generator is $\G=\B\A$.

\paragraph{Wrapped binary convolution.}
Fix a memory $m\geq1$.  Initialize $y_j:=0$ for $j\leq0$.  For each time
$t$, sample independent bits $b_{t,1},\ldots,b_{t,m-1}$ and compute
\begin{equation}
 y_t:=u_t+y_{t-m}+\sum_{j=1}^{m-1}b_{t,j}y_{t-j}
 \quad\text{over }\F_2.
 \label{eq:wrapped-convolution}
\end{equation}
The coefficient of the oldest state bit is fixed to one.  This is the wrapped
EC convolution of~\cite{EC}.  Its coefficient on $u_t$ is one, so every
sampled convolution is invertible.

\paragraph{Finite-field convolution.}
Fix a prime power $p$, choose the field $\F_p$, and fix memory $m$.  At time
$t$, independently sample
$\boldsymbol\alpha_t\in\F_p^m$ and compute
\begin{equation}
 y_t:=u_t+
 \left\langle\boldsymbol\alpha_t,
 (y_{t-1},\ldots,y_{t-m})\right\rangle.
 \label{eq:prime-convolution}
\end{equation}
This map is also invertible.  Sampling from the full field is essential for
the reduced-state analysis in Section~\ref{sec:prime-fields}.

\subsection{Expander distributions}

We use three principal distributions for $\B$.

\paragraph{Bernoulli incidence.}
Over $\F_2$, independently set each entry of $\B$ to one with probability
$\rho$ and to zero otherwise.  We write
$\B\gets\mathcal B^{\mathrm{Ber}}_{k,n,\rho}$.  A row has expected weight
$\rho n$.

Over $\F_p$, an entry is zero with probability $1-\rho$.  Otherwise it is
sampled from $\F_p^*$.  We call the nonzero value an edge label.  All
incidence and label choices are independent unless stated otherwise.

\paragraph{Region-stratified left regularity.}
Assume $n=d_L\ell$ and partition the right coordinates into $d_L$ consecutive
regions of length $\ell$.  For each left vertex and each region, choose one
right coordinate uniformly and independently.  Place one edge there.  Over
$\F_p$, independently label the edge with an element of $\F_p^*$.

Every left vertex has degree $d_L$, with one edge in every region.  Right
degrees remain random.  We call this distribution a
\emph{region-stratified left-regular expander}.  We use
\emph{regular expander} when the one-sided meaning is already clear.

\paragraph{Two-sided regularity.}
Assume $k=d_R\ell$ and $n=d_L\ell$, so $kd_L=nd_R$.  Within each region,
sample a permutation of the $k$ left vertices.  Split the permuted list into
$\ell$ groups of size $d_R$ and connect each group to one right coordinate.
Every left vertex has one edge per region, and every right vertex has degree
$d_R$.  We call this a \emph{two-sided regular expander}.

At rate one half, $d_R=d_L/2$.  The finite-field EC certificates in
Section~\ref{sec:certificates} use this distribution.  The binary regular
EA certificate uses left regularity.  The binary EC frontier includes both
left-regular and two-sided regular ensembles.

\begin{table}[t]
\centering
\small
\caption{Random choices in the principal ensembles.}
\label{tab:ensemble-definitions}
\begin{tabular}{@{}llll@{}}
\toprule
name & field & expander & recursive map \\
\midrule
binary EA & $\F_2$ & Bernoulli & fixed accumulator \\
binary wrapped EC & $\F_2$ & Bernoulli & independent wrapped convolution \\
regular binary EA & $\F_2$ & left regular & fixed accumulator \\
regular binary EC & $\F_2$ & left regular & independent wrapped convolution \\
labeled $p$-ary EA & $\F_p$ & labeled Bernoulli or regular & fixed accumulator \\
labeled $p$-ary EC & $\F_p$ & labeled two-sided regular & independent full-field convolution \\
\bottomrule
\end{tabular}
\end{table}

The word \emph{independent} in Table~\ref{tab:ensemble-definitions} means that
the expander and recursive-map randomness are sampled independently.  Within
the regular expander, incidences have the dependencies imposed by the stated
degree constraints.

\paragraph{Sampling and cost.}
The following procedures make the setup cost explicit.  For a left-regular
expander, process regions $j=1,\ldots,d_L$.  For every left vertex $i$, sample
$c_{i,j}\gets\{1,\ldots,\ell\}$ and connect $i$ to coordinate
$(j-1)\ell+c_{i,j}$.  For a two-sided regular expander, sample a permutation
$\pi_j$ of the $k$ left vertices in each region.  Connect
$\pi_j((q-1)d_R+1),\ldots,\pi_j(qd_R)$ to the $q$th coordinate of that region.
Over $\F_p$, sample one independent label from $\F_p^*$ for every generated
edge.

Both procedures generate exactly $kd_L=nd_R$ edges.  A materialized encoder
stores that many indices and, over $\F_p$, the same number of labels.  It uses
one scatter-add per binary edge or one multiply-add per labeled edge.
Left-regular setup samples $kd_L$ independent indices.  Two-sided setup uses
$d_L$ permutations of length $k$; a Fisher--Yates implementation takes
$O(kd_L)$ random words and $O(k)$ reusable workspace.  A streaming encoder
may regenerate one region at a time from a public seed and reduce storage to
$O(k)$ indices.  The certificate comparisons count edges only.  They do not
claim that permutations, transposition, and field multiplications have equal
cost.
